\section{Entscheidungen}

\subsection{Die Umgebung \texttt{decision}}

Die Umgebung \verb|decision| bildet einen vollständigen Entscheidungsblock
mit einer Bedingung und mehreren möglichen Ausgängen.

Ein gültiger Block enthält immer:

\begin{enumerate}
\item genau eine \verb|condition|,
\item mindestens eine \verb|pass|,
\item mindestens eine \verb|fail|.
\end{enumerate}

Werden diese Strukturregeln verletzt, bricht die Kompilierung mit einer
verständlichen Fehlermeldung ab.

\subsubsection*{Syntax}

\begin{Verbatim}
\begin{decision}

    \begin{condition}
    Bedingung
    \end{condition}

    \begin{pass}[Optionaler Titel]
    Erfolgstext
    \end{pass}

    \begin{fail}[Optionaler Titel]
    Misserfolgstext
    \end{fail}

\end{decision}
\end{Verbatim}

\subsubsection*{Hinweise}

\begin{itemize}
\item Text direkt in \verb|decision| außerhalb der Kind-Umgebungen wird ignoriert.
\item Leerräume und Kommentare zwischen Kind-Umgebungen werden ignoriert.
\item Die Reihenfolge im Quelltext beeinflusst die Ausgabe nicht:
      dargestellt wird immer \emph{condition}, dann \emph{pass}, dann \emph{fail}.
\item Die Reihenfolge innerhalb der \verb|pass|-Einträge bzw. \verb|fail|-Einträge
      bleibt erhalten.
\end{itemize}

\subsubsection*{Beispiel}

\paragraph{Quellcode}

\begin{Verbatim}
\begin{decision}

    \begin{condition}
    Wahrnehmungsprobe -2
    \end{condition}

    \begin{pass}[0]
    Du siehst ein seltsames Glühen.
    \end{pass}

    \begin{pass}[Kritischer Erfolg]
    Du erkennst sofort die Energiequelle eines elektronischen Attentäters.
    \end{pass}

    \begin{fail}
    Der Korridor ist dunkel und wirkt unheimlich.
    \end{fail}

    \begin{fail}[Kritischer Fehlschlag]
    Für dich sieht der Gang vollkommen harmlos aus.
    \end{fail}

\end{decision}
\end{Verbatim}

\paragraph{Ergebnis}

\begin{decision}

    \begin{condition}
    Wahrnehmungsprobe -2
    \end{condition}

    \begin{pass}[0]
    Du siehst ein seltsames Glühen.
    \end{pass}

    \begin{pass}[Kritischer Erfolg]
    Du erkennst sofort die Energiequelle eines elektronischen Attentäters.
    \end{pass}

    \begin{fail}
    Der Korridor ist dunkel und wirkt unheimlich.
    \end{fail}

    \begin{fail}[Kritischer Fehlschlag]
    Für dich sieht der Gang vollkommen harmlos aus.
    \end{fail}

\end{decision}


\subsection{Die Umgebung \texttt{condition}}

\verb|condition| definiert die Bedingung des Entscheidungsblocks
(oberer, bläulicher Bereich).

Sie darf nur \emph{direkt} innerhalb einer \verb|decision|-Umgebung stehen
und genau einmal vorkommen.

\subsubsection*{Syntax}

\begin{Verbatim}
\begin{condition}
Frage oder Bedingung
\end{condition}
\end{Verbatim}

\subsubsection*{Beispiel}

\paragraph{Quellcode}

\begin{Verbatim}
\begin{decision}
    \begin{condition}
    Ist die Tür verschlossen?
    \end{condition}
    \begin{pass}
    Nein, sie ist offen.
    \end{pass}
    \begin{fail}
    Ja, du brauchst einen Dietrich.
    \end{fail}
\end{decision}
\end{Verbatim}

\paragraph{Ergebnis}

\begin{decision}
    \begin{condition}
    Ist die Tür verschlossen?
    \end{condition}
    \begin{pass}
    Nein, sie ist offen.
    \end{pass}
    \begin{fail}
    Ja, du brauchst einen Dietrich.
    \end{fail}
\end{decision}


\subsection{Die Umgebung \texttt{pass}}

\verb|pass| definiert einen möglichen Erfolgs-Eintrag
(grünlicher Bereich).

Es können beliebig viele \verb|pass|-Einträge angegeben werden,
mindestens jedoch einer.

\subsubsection*{Syntax}

\begin{Verbatim}
\begin{pass}[Optionaler Titel]
Erfolgseintrag
\end{pass}
\end{Verbatim}

\subsubsection*{Titel und Ausrichtung}

\begin{itemize}
\item Einträge dürfen einen optionalen Titel in eckigen Klammern besitzen.
\item Bei genau einem Eintrag wird der Titel (falls vorhanden) fett über dem Inhalt gesetzt.
\item Bei mehreren Einträgen werden Titel links neben dem Inhalt ausgerichtet;
      die Einrückung richtet sich nach der breitesten \verb|pass|-Überschrift.
\item Einträge ohne Titel gelten als leerer Titel und reservieren denselben Titelbereich.
\end{itemize}

\subsubsection*{Beispiel}

\paragraph{Quellcode}

\begin{Verbatim}
\begin{decision}
    \begin{condition}
    Probe geschafft?
    \end{condition}
    \begin{pass}[Standard]
    Du findest eine verwertbare Spur.
    \end{pass}
    \begin{pass}[Kritischer Erfolg]
    Du findest die Spur und gewinnst sofort einen taktischen Vorteil.
    \end{pass}
    \begin{fail}
    Du findest nichts.
    \end{fail}
\end{decision}
\end{Verbatim}


\subsection{Die Umgebung \texttt{fail}}

\verb|fail| definiert einen möglichen Misserfolgs-Eintrag
(rötlicher Bereich).

Es können beliebig viele \verb|fail|-Einträge angegeben werden,
mindestens jedoch einer.

\subsubsection*{Syntax}

\begin{Verbatim}
\begin{fail}[Optionaler Titel]
Misserfolgseintrag
\end{fail}
\end{Verbatim}

\subsubsection*{Titel und Ausrichtung}

Die Layoutregeln entsprechen \verb|pass|, werden aber unabhängig berechnet:

\begin{itemize}
\item Für die Einrückung werden nur \verb|fail|-Titel betrachtet.
\item \verb|pass|-Titel beeinflussen die \verb|fail|-Ausrichtung nicht.
\end{itemize}

\subsubsection*{Beispiel}

\paragraph{Quellcode}

\begin{Verbatim}
\begin{decision}
    \begin{condition}
    Schleichprobe gelungen?
    \end{condition}
    \begin{pass}
    Ihr bleibt unentdeckt.
    \end{pass}
    \begin{fail}[Fehlschlag]
    Eine Wache wird misstrauisch.
    \end{fail}
    \begin{fail}[Kritischer Fehlschlag]
    Alarm wird ausgelöst.
    \end{fail}
\end{decision}
\end{Verbatim}


\subsection{Verschachtelung und Gültigkeit}

Kind-Umgebungen dürfen nicht ineinander verschachtelt werden.

Ungültig sind insbesondere:

\begin{itemize}
\item \verb|condition| innerhalb von \verb|condition|, \verb|pass| oder \verb|fail|,
\item \verb|pass| innerhalb einer anderen Kind-Umgebung,
\item \verb|fail| innerhalb einer anderen Kind-Umgebung.
\end{itemize}

Kind-Umgebungen sind nur direkt unter \verb|decision| erlaubt.
Bei Verstößen bricht die Kompilierung mit einer Fehlermeldung ab.


\subsection{Leere Einträge}

Leere \verb|pass|- und \verb|fail|-Einträge sind gültig.

\subsubsection*{Beispiele}

\begin{Verbatim}
\begin{pass}
\end{pass}

\begin{fail}[Kritischer Fehlschlag]
\end{fail}
\end{Verbatim}


\subsection{Darstellung und Seitenumbrüche}

Der Entscheidungsblock wird als zusammenhängende Einheit dargestellt:

\begin{itemize}
\item Breite: 90\% der verfügbaren Zeilenbreite,
\item horizontal zentriert,
\item Inhalt linksbündig,
\item farblich getrennte Bereiche für \verb|condition|, \verb|pass| und \verb|fail|,
\item deutliche Trennlinien zwischen den Bereichen.
\end{itemize}

Seitenumbrüche innerhalb langer Inhalte sind erlaubt.
Inhalte werden über Seiten hinweg fortgesetzt und nicht abgeschnitten.
